% ===========================================================================
%  Bagian 7 -- Utilitas (benchmark, diagnose, compare, export, map, simulasi)
% ===========================================================================
\section{Lapisan utilitas: benchmarking, diagnostik, perbandingan, ekspor, peta}\label{sec:utilitas}

Kelima kelompok fungsi di bawah ini bekerja pada objek apa pun berkelas
\code{fastsae} (atau turunannya) dan membentuk ``alur kerja'' setelah sebuah
model selesai dipasang pada data: \emph{kalibrasi} $\to$ \emph{diagnosis}
$\to$ \emph{perbandingan} $\to$ \emph{pelaporan/pemetaan}. Semua
rumus adalah rumus persis yang dikodekan; selisih terhadap rujukan ditandai
\emph{implementation note} dan dikumpulkan di \code{DISCREPANCIES.md}
bagian~G (ID \texttt{UT-1}--\texttt{UT-16}).

\subsection{Benchmarking: \texttt{benchmark\_sae}}\label{sec:bm-sae}

\paragraph{Definisi.} Agregat per grup dihitung dari bobot yang sudah
disterilkan oleh \code{type}:
\begin{equation}
  w^{\text{calc}}_{ig}=
  \begin{cases}
    w_{ig}/\sum_{k\in g}w_{ik} & \text{\code{type = "mean"}}\\
    w_{ig} & \text{\code{type = "total"}}
  \end{cases}
  \qquad
  A_g=\sum_{i\in g}w^{\text{calc}}_{ig}\,y_i,
  \label{eq:bm-agg}
\end{equation}
dan target $T_g$ dipaksa tercapai lewat empat metode
(\code{R/benchmark.R:451-512}):
\begin{align}
  \text{\textbf{ratio:}}\quad
    &\lambda_g=\frac{T_g}{A_g},\qquad \tilde y_i=y_i\,\lambda_g;
    \label{eq:bm-ratio}\\
  \text{\textbf{difference:}}\quad
    &\delta_g=T_g-A_g,\qquad
      \tilde y_i=y_i+\frac{\delta_g}{\sum_i w^{\text{calc}}_{ig}};
    \label{eq:bm-diff}\\
  \text{\textbf{optimal:}}\quad
    &\lambda_g=\frac{\delta_g}{\sum_i \tilde m_i},\qquad
      \tilde y_i=y_i+\frac{\tilde m_i}{w^{\text{calc}}_{ig}}\,\lambda_g;
    \label{eq:bm-opt}\\
  \text{\textbf{logit:}}\quad
    &\text{cari }\alpha\text{ dari }\textstyle
      \sum_i w^{\text{calc}}_{ig}\,
      \mathrm{logit}^{-1}\bigl(\mathrm{logit}(y_i)+\alpha\bigr)=T_g,\notag\\
    &\tilde y_i=\mathrm{logit}^{-1}\bigl(\mathrm{logit}(y_i)+\hat\alpha\bigr),
    \label{eq:bm-logit}
\end{align}
dengan $\tilde m_i=\max(\mathrm{mse}_i,10^{-8})$; galat
\code{ratio} dicatat bila $\abs{A_g}<$ \code{double.eps}. Persamaan
\eqref{eq:bm-logit} dicari lewat \code{uniroot} pada bracket
$[-20,20]$, diperluas ke $[-50,50]$ bila tandanya sama, dengan
\code{tol = 1e-9}; persyaratan $0<y<1$ (dan $0<T_g<1$ untuk
\code{type = "mean"}) diperiksa lebih dulu. Kepatuhan constraint
$\sum_i w^{\text{calc}}_{ig}\tilde y_i=T_g$ terverifikasi aljabar untuk
keempat metode.

\implnote{Metode \code{optimal} memakai alokasi
$\delta_i=(\mathrm{mse}_i/w_i)\lambda$ sementara solusi KKT standar dari
$\min\sum\delta_i^2/\mathrm{mse}_i$ s.t. $\sum w_i\delta_i=\Delta$ adalah
$\delta_i\propto w_i\mathrm{mse}_i$. Constraint tetap tercapai, tetapi
alokasinya berbeda ketika bobot tidak seragam (UT-6); variabelnya juga
dinamai \code{inv\_prec} padahal isinya \code{mse}.}

\paragraph{Dua tahap hierarkis.} Aktif bila \code{national\_target} diberikan
dan ada $>1$ grup (\code{R/benchmark.R:318}): STAGE~1 menyetel target grup
$T_g$ agar agregatnya menyamai target nasional (empat metode yang sama,
\code{W\_share}=porsi bobot grup), STAGE~2 menyetel estimasi level kecil ke
$T_g$; keduanya dicatat \code{stage\_1\_summary} dan diverifikasi ulang.
\begin{algorithm}[htbp]
\caption{\texttt{benchmark\_sae} --- alur sebenarnya (\texttt{R/benchmark.R:120-573})}\label{alg:bm}
\KwIn{objek model (atau vektor numerik), \code{target}, \code{weight},
  \code{method} $\in$\{ratio, difference, optimal, logit\}, \code{group},
  \code{type}, \code{national\_target}, \code{outer\_target}, \code{mse}}
\KwOut{data frame kelas \code{fastsae\_benchmark} + atribut verifikasi}
estimasi $\leftarrow$ \texttt{df\_hb \%||\% df\_ebp \%||\% df\_eblup};
  $y\leftarrow$ \code{hb/ebp/eblup/est}; $\mathrm{mse}\leftarrow$ \code{mse} atau $\mathrm{sd}^2$ (142-153)\;
bobot $\leftarrow$ nama kolom / vektor / seragam $1$ (157-173)\;
grup $\leftarrow$ nama kolom / vektor / \code{"All"} (178-191)\;
$A_g,T_g\leftarrow$ \eqref{eq:bm-agg} + parsing target (328-393)\;
\If{\code{national\_target} $\ne$ NULL \textbf{dan} $>1$ grup}{STAGE 1: setel $T_g$ agar $\sum_g W_{\text{share},g}T_g=T_{\text{nas}}$ (395-438)\;}
STAGE 2: $\tilde y\leftarrow$ \eqref{eq:bm-ratio}--\eqref{eq:bm-logit} per grup (440-512)\;
verifikasi $|\sum w^{\text{calc}}\tilde y-T_g|$; status \code{CONSISTENT} bila $<10^{-5}$ (534-560)\;
susun \code{domain, weight, original, benchmarked, adjustment, rel\_adjustment\_pct, target} (515-531)\;
\end{algorithm}

\paragraph{Keluaran.} \code{data.frame} kelas \code{c("fastsae\_benchmark","data.frame")}
dengan atribut \code{method}, \code{type}, \code{hierarchical},
\code{stage1\_summary}, \code{verification}, \code{national\_verification};
metode S3: \code{print} (status konsistensi per grup), \code{summary},
\code{autoplot} (sebaran \code{original} vs \code{benchmarked} dengan garis 1:1),
dan \code{plot} yang bila diberi \code{sf\_geom} meneruskan ke \fct{map\_sae}.
Rujukan yang memverifikasi pendekatan restrisi/benchmarking di tingkat
estimasi: \citet{rao2015small} (bab 10), \citet{you2002pseudo} dan
\citet{wang2008small}.

\subsection{Diagnostik: \texttt{diagnose}}\label{sec:diagnose}

\code{diagnose()} adalah fungsi biasa (bukan generic S3) dan mensyaratkan
objek \code{fastsae} (\code{R/diagnose.R:53-60}). Kuantitas dasar
(\code{R/diagnose.R:84-130}):
\begin{equation}
  \mathrm{RSE}^{\text{direct}}_d=100\frac{\sqrt{\psi_d}}{\abs{y_d}},
  \qquad
  \mathrm{eff}_d=\frac{\psi_d}{\mathrm{mse}_d},
  \qquad
  r_d=\frac{y_d-\hat\theta_d}{\sqrt{\psi_d}},
  \qquad
  \widehat{\Pr}(\mathrm{RSE}<T)=\frac{\#\{\mathrm{RSE}<T\}}{n},
  \label{eq:diag-basic}
\end{equation}
dengan ambang default $T=25$\,\% dan area tanpa sampel dihitung
terpisahnya (\code{N\_sampled}/\code{N\_unsampled}).

\paragraph{Kalibrasi Brown et al.\ (2001) (asal: kode).} Untuk
minimal lima pasang $(y,\hat\theta)$ ter-sampel:
\begin{align}
  y_d&=\alpha+\beta\hat\theta_d+e_d \quad\text{(OLS polos)}, \label{eq:brown-reg}\\
  F&=\frac{c'\,V^{-1}c}{2},\quad c=(\hat\alpha,\hat\beta)'-(0,1)',
    \quad c\sim F_{2,\,n-2}, \label{eq:brown-wald}\\
  W&=\sum_d\frac{(y_d-\hat\theta_d)^2}{\psi_d+\mathrm{mse}_d},
    \quad W\sim\chi^2_{n_w} \text{ (asimptotik)}. \label{eq:brown-gof}
\end{align}
\code{is\_unbiased} diputuskan dari $p$-value $F$; \code{is\_good\_fit}
diputuskan dari penerimaan interval dua sisi
$\chi^2_{\alpha/2,n_w}\le W\le\chi^2_{1-\alpha/2,n_w}$.

\implnote{Empat hal pada paragraf di atas adalah temuan lapangan dan dicatat
sebagai UT-1--UT-3: regresi memakai OLS tanpa bobot (bukan WLS
$\propto1/(\psi+\mathrm{mse})$ seperti yang lazim dikutip), statistik Wald
dibagi $2$ lalu dibandingkan $F(2,n-2)$ (bukan $\chi^2(2)$), keputusan GOF
memakai interval dua sisi sedangkan $p$-value yang dicetak ekor-atas tunggal
sehingga status bisa \code{CAUTION} walau $p\ge\alpha$, dan tidak ada teks
rujukan Brown et al.\ dalam repo --- dokumen resminya laporan internal ONS
yang tidak dapat diverifikasi di Crossref (lihat \code{TODO\_REFERENCES.md}).}

\paragraph{Moran's I residual} (\code{R/diagnose.R:193-246}) dihitung pada
residual mentah area ter-sampel (bukan terstandarisasi), $W$ dipotong ke
baris ter-sampel, minimal lima area:
\begin{align}
  I&=\frac{n}{S_0}\cdot\frac{\sum_{ij}w_{ij}z_iz_j}{\sum_i z_i^2},
    \label{eq:moran-I}\\
  \Var(I)&=\frac{
    n\bigl[(n^2-3n+3)S_1-nS_2+3S_0^2\bigr]
    -\kappa\bigl[(n^2-n)S_1-2nS_2+6S_0^2\bigr]}
    {(n-1)(n-2)(n-3)S_0^2}-\bigl(\mathbb E I\bigr)^2,
  \label{eq:moran-var}
\end{align}
dengan $S_0=\sum_{ij}w_{ij}$, $z=r-\bar r$,
$S_1=\tfrac12\sum(w_{ij}+w_{ji})^2$, $S_2=\sum_i(\sum_jw_{ij}+\sum_jw_{ji})^2$,
$\kappa=n\sum z^4/(\sum z^2)^2$, $\mathbb E I=-1/(n-1)$, dan
$\Var\ge10^{-8}$; $z$-statistik dua sisi $2\Phi(-\abs{z})$.

\paragraph{Metrik simulasi dan diagnostik Bayes.} Bila \code{truth} diberikan:
$\mathrm{rb}=100(\hat\theta-\theta)/\theta$, $\mathrm{rrmse}=100\abs{\hat\theta-\theta}/\abs{\theta}$
(dinamakan ``Mean Relative RMSE'' padahal ini rerata galat relatif absolut,
UT-4), dan \code{coverage} dari kolom CI apa adanya (UT-5). Untuk model
Bayes, DIC/WAIC/marginal likelihood diambil dari \code{goodness}, CPO/PIT dari
\code{fit\$cpo}, jumlah \code{failure}, lalu uji KS
\code{ks.test(pit, "punif")} (syarat lima nilai PIT di $(0,1)$).

\paragraph{Status akhir} (\code{R/diagnose.R:334-349}) berurutan:
\code{PASS} bila semua tes lolos, lalu \code{CAUTION} untuk bias, \code{CAUTION}
untuk GOF, lalu \code{WARNING} bila proporsi RSE$\ge$ ambang terlalu besar.
Bila suatu tes tidak tersedia dianggap lolos.

\begin{algorithm}[htbp]
\caption{\texttt{diagnose} --- alur sebenarnya (\texttt{R/diagnose.R:53-386})}\label{alg:diagnose}
\KwIn{objek \code{fastsae}, \code{W}, \code{truth}, \code{rse\_threshold} $=25$, \code{alpha\_level} $=0.05$}
\KwOut{list kelas \code{fastsae\_diagnose} + string \code{status}}
ekstrak \code{df}, $y$, $\hat\theta$, $\psi$, $\mathrm{mse}$ (63-73)\;
\eqref{eq:diag-basic}: RSE direct/SAE, efisiensi, residual, proporsi andal (84-130)\;
\If{$n_{\text{pasang}}\ge5$}{regresi \eqref{eq:brown-reg}; Wald \eqref{eq:brown-wald}; GOF \eqref{eq:brown-gof} (133-191)}\;
\If{\code{W} tersedia \textbf{dan} $n\ge5$}{Moran's I \eqref{eq:moran-I}--\eqref{eq:moran-var} pada residual terpotong (193-246)}\;
\If{\code{truth} diberikan}{rb, rrmse, coverage (251-271)}\;
\If{objek Bayes}{DIC/WAIC/CPO/PIT + uji KS (276-332)}\;
\code{status} $\leftarrow$ aturan berurutan pada teks (341-349)\;
susun \code{df\_diag} (\code{domain, direct\_y, sae\_pred, vardir, mse, direct\_rse, sae\_rse, eff\_ratio, residuals, std\_residuals, is\_sampled}[, \code{cpo, pit}]) (351-382)\;
\end{algorithm}

Keluaran dicetak \code{print.fastsae\_diagnose} dalam lima bagian (efisiensi,
kalibrasi Brown, Moran's I, simulasi, Bayes) plus penilaian akhir; grafik
lewat \code{autoplot} dengan tipe \code{all}, \code{calibration},
\code{rse}, \code{residuals}, \code{qq}, \code{pit}, \code{cpo}.
Tidak ada \code{summary} maupun \code{plot} untuk kelas ini (UT-13).

\subsection{Perbandingan dua model: \texttt{compare\_sae}}\label{sec:compare}

\code{compare\_sae(model1, model2, names, thresholds = c(20,30))} hanya
memiliki method \code{.default} (UT-14); kunci gabung adalah
\code{domain}, atau \code{c(domain, subarea)}, atau \code{c(domain, time)}
(inner join; error bila tumpang tindih nol). Metrik yang dihitung
(\code{R/compare\_sae.R:112-176}):
\begin{equation}
  r=\operatorname{cor}(\hat\theta_1,\hat\theta_2),\quad
  \rho=\text{Spearman},\quad
  \mathrm{MAE}=\overline{\abs{d}},\quad
  \mathrm{RMSD}=\sqrt{\overline{d^{2}}},\quad
  d=\hat\theta_2-\hat\theta_1,
  \label{eq:compare-metrics}
\end{equation}
plus rasio $\mathrm{mse}_1/\mathrm{mse}_2$, jumlah domain dengan
$\mathrm{mse}_2<\mathrm{mse}_1$, dan $\Delta\overline{\mathrm{RSE}}$.
\emph{Tidak ada} argumen \code{truth}, sehingga RMSD/MAE di sini adalah
selisih antarmodel, bukan galat terhadap kebenaran.
\code{thresholds} hanya dipakai sebagai garis horizontal pada plot
\code{rse} (UT-14). Lima tipe \code{autoplot}: \code{scatter} (garis 1:1),
\code{comparison} (default, dodged), \code{difference} (lollipop),
\code{mse} (bar dodged), \code{rse} (garis + ambang).

\subsection{Flag reliabilitas dan ekspor: \texttt{export\_sae}}\label{sec:export}

\code{add\_reliability\_flags} memotong RSE dengan
$\texttt{cut}(\text{breaks}=(-\infty,20,30,\infty))$ menjadi tiga tingkat
\emph{Reliable (\textless 20\%)}, \emph{Use with Caution (20--30\%)},
\emph{Unreliable ($\ge$30\%)}; bila kolom RSE tidak ada, ia menghitung
$100\sqrt{\mathrm{mse}}/\abs{\hat\theta}$. \implnote{Dokumen menyebut fungsi
menerima ``objek atau data frame'' tetapi nilai kembalian selalu berupa
data frame saja --- penugasan balik ke \code{object\$df\_*} tidak efektif
karena aturan copy-on-modify R (UT-9).}

\code{export\_sae(object, file, benchmark, thresholds, overwrite)} hanya
menerima ekstensi \code{xlsx}/\code{csv}; sheet yang ditulis:
\code{Estimates}, \code{Model\_Summary}, \code{Benchmarked} dan
\code{Benchmarked\_Groups} (bila hierarkis) --- bisa empat, bukan ``hingga
tiga'' seperti dokumentasi (UT-10). Penulisan xlsx memakai
\pkg{writexl}, fallback \pkg{openxlsx}, fallback terakhir CSV berisi
\code{Estimates} saja; sheet \code{Benchmarked} tidak pernah ikut pada mode
csv. Sheet \code{Estimates} memuat
$\mathrm{SE}=\sqrt{\mathrm{mse}}$, $\mathrm{CI}_{95\%}=\hat\theta\pm1.96\,\mathrm{SE}$
(bila kolom CI tidak ada), serta flag reliabilitas; sheet
\code{Model\_Summary} berupa pasangan \code{Parameter/Value} (paket, tipe
model, method, konvergensi, komponen ragam, goodness of fit, koefisien
$+$ SE $+$ $p$).

\subsection{Peta choropleth: \texttt{map\_sae}}\label{sec:map}

Empat method: \code{.fastsae}, \code{.fastsae\_benchmark}, \code{.list},
\code{.default}; semuanya mensyaratkan \pkg{sf} dan menghasilkan
\code{ggplot}. Lima tipe \code{type}:
\begin{itemize}
  \item \code{estimate} --- \code{geom\_sf} + viridis;
  \item \code{rse} --- \code{option = "magma"}, arah terbalik;
  \item \code{reliability} --- tiga warna manual
        \texttt{\#198754} / \texttt{\#ffc107} / \texttt{\#dc3545}
        sesuai tingkat \S\ref{sec:export};
  \item \code{comparison} --- facet Direct Survey vs SAE (otomatis jatuh ke
        \code{estimate} bila kolom direct seluruhnya \code{NA});
  \item \code{difference} --- $\hat\theta-y$ dengan \code{gradient2}
        \texttt{\#2b83ba}/\texttt{\#ffffbf}/\texttt{\#d7191c}.
\end{itemize}
Pencocokan kunci otomatis memakai daftar kandidat
\code{domain, area, id, code, kd\_kab, kode, kabupaten, provinsi, district,
region} dan mengadopsi pasangan dengan tumpang tindih terbesar
\emph{bila mencapai} $40\%$ jumlah domain; gagal $\to$ error.

\implnote{Argumen \code{thresholds = c(20,30)} diteruskan tetapi tidak pernah
dipakai: pemotongan reliabilitas hard-coded di extractor (UT-12). Pemetaan
dua model menyusun \code{data.frame(matched\_sf[, c("domain","geometry")],\dots)},
yang akan gagal bila kolom geometri \code{sf} bukan \code{"geometry"}
(UT-16).}

\subsection{Metode S3 dan grafik}\label{sec:s3}

Tabel~\ref{tab:s3} merangkum metode yang terdaftar di \code{NAMESPACE:23-52}.

\begin{table}[htbp]
\centering
\caption{Metode S3 utama beserta file sumbernya.}
\label{tab:s3}
\small
\begin{tabularx}{\textwidth}{@{}l l X@{}}
\toprule
Generic & Objek & Perilaku ringkas \\
\midrule
\code{print} & \code{fastsae} & Call, konvergensi + \code{n\_iter},
  tipe model, komponen ragam, \code{printCoefmat(estcoef)},
  \code{head(df,6)} \\
\code{summary} & \code{fastsae} & list 20 komponen; dicetak
  \code{print.summary.fastsae} \\
\code{coef}, \code{fitted}, \code{residuals} & \code{fastsae} &
  \code{estcoef\$beta}; \code{hb/ebp/eblup}; $y-\hat\theta$ \\
\code{plot}, \code{autoplot} & \code{fastsae} & 7 tipe: \code{comparison},
  \code{ribbon}, \code{mse}, \code{estimates}, \code{scatter}, \code{rse},
  \code{map} \\
\code{autoplot} & \code{list fastsae} & \code{comparison}, \code{mse},
  \code{scatter} (butuh \code{names(x)} --- UT-15) \\
\code{autoplot} & \code{fastsae\_diagnose} & \code{all},
  \code{calibration}, \code{rse}, \code{residuals}, \code{qq},
  \code{pit}, \code{cpo} \\
\code{print}, \code{summary}, \code{autoplot}, \code{plot} &
  \code{fastsae\_benchmark} & lihat \S\ref{sec:bm-sae} \\
\code{print}, \code{summary}, \code{autoplot}, \code{plot} &
  \code{fastsae\_comparison} & lihat \S\ref{sec:compare} \\
\bottomrule
\end{tabularx}
\end{table}

Pada \code{autoplot} model tunggal, bila kolom CI belum ada, pita dibuat
$\hat\theta\pm1.96\sqrt{\mathrm{mse}}$; tipe \code{rse} membandingkan RSE
SAE vs RSE direct dengan dua ambang \code{c(20,30)}.

\subsection{Data simulasi}\label{sec:sim}

Tiga pembangkit tersedia, semuanya menghasilkan matriks ketetanggaan lewat
\fct{sim\_spatial\_weights} (tipe \code{knn}/\code{grid}/\code{ring}, gaya
\code{B} biner atau \code{W} row-standar; kNN disimetrisasi
$A\leftarrow\max(A,A')$).

\paragraph{DGP area-level} (\code{sim\_area\_data}, \code{R/sim\_area\_data.R:63-166}):
\begin{align}
  v&=(\Imat-\rho W_{\text{std}})^{-1}\varepsilon,\quad
    \tilde v=\text{standarisasi}(v),\quad
    \tilde u=\text{standarisasi}(N(0,1)),\label{eq:sim-spatial}\\
  u&=\sigma_u\bigl(\sqrt{\varphi}\,\tilde v+\sqrt{1-\varphi}\,\tilde u\bigr),
    \qquad
    \eta=x'\beta+u, \label{eq:sim-bym2}\\
  y\mid\eta&\sim
    \begin{cases}
      \eta+\mathcal N(0,\psi) & \psi\sim U(0.04,\,0.16)\\
      \text{Poisson}\bigl(\text{exposure}\cdot
        e^{\operatorname{clip}(\eta/2-0.5,\,-4,\,4)}\bigr) & \text{offset}\\
      \text{Binomial}\bigl(\text{trials},\ \mathrm{logit}^{-1}(\eta-1)\bigr) & \text{peluang}\\
      \text{NB}\bigl(\theta=3,\ \mu=15\,e^{\operatorname{clip}(\eta/2,\,-3,\,4)}\bigr) &
        \text{gamma-Poisson}\\
      \text{Beta}\bigl(\mu\kappa,\ (1-\mu)\kappa\bigr) & \kappa=25\\
      \text{Gamma}\bigl(k=5,\ \mu/k\bigr) & k=5
    \end{cases}
    \label{eq:sim-resp}
\end{align}
dengan $x_1\sim N(2,1)$, $x_2\sim U(0,5)$, parameter dipotong
$\rho\in[-0.95,0.95]$, $\varphi\in[0,1]$, dan \code{n\_unsampled} domain
diset \code{NA}; \code{exposure} $\sim\mathrm{round}\,U(80,400)$ dan
\code{trials} $\sim\mathrm{round}\,U(50,250)$. Rata-rata pada Beta adalah
$\mu=\operatorname{clip}\bigl(\mathrm{logit}^{-1}(\eta-1),\,0.02,\,0.98\bigr)$
(dan $y$ dipotong ke $[0.001,0.999]$), pada Gamma
$\mu=\exp\bigl(\operatorname{clip}(\eta/2,\,-3,\,3)\bigr)$, dan pada NB
$\mu=15\exp\bigl(\operatorname{clip}(\eta/2,\,-3,\,4)\bigr)$. \eqref{eq:sim-bym2} adalah campuran gaya BYM2 yang sama
dengan \eqref{eq:bym2}.

\paragraph{DGP panel} (\code{sim\_series\_data}) menambah efek spasial
$u_1=\sigma_s\,\text{standar}((\Imat-\rho_sW_{\text{std}})^{-1}\varepsilon)$,
proses AR(1) per domain $u_2[d,t]=\rho_t u_2[d,t-1]+\mathcal N(0,\sigma_t\sqrt{1-\rho_t^2})$
(varian stasioner $\sigma_t^2$), dan tren deterministik
(\code{linear}/\code{random\_walk}/\code{ar1}/\code{none}); pengurutan baris
\code{domain-major} memenuhi syarat \fct{eblup\_stfh}.

\implnote{Dataset terlampir berbeda dari keluaran pembangkit: \code{sim\_area}
memakai kolom \code{area} dengan 14 kolom tanpa \code{truth\_gaussian},
sedangkan \code{sim\_area\_data()} menghasilkan 15 kolom dengan
\code{domain} dan \code{truth\_gaussian} (UT-11); hal yang sama berlaku
untuk \code{sim\_panel}.}

\subsection{Rujukan lapisan utilitas}\label{sec:util-ref}

Blok roxygen \code{R/benchmark.R:79-93} memuat lima rujukan dan
\code{R/diagnose.R:27-33} memuat tiga; yang berhasil diverifikasi dipakai di
sini: \citet{rao2015small} (bab benchmarking), \citet{you2002pseudo},
\citet{datta2010bayesian} (benchmarking Bayes), dan \citet{berg2014small}.
Rujukan yang tidak dapat diverifikasi (Steorts et al.\
2014; Brown et al.\ 2001; Cliff \& Ord 1981) tidak dikutip sebagai
referensi bibliografis dan didaftarkan di \code{TODO\_REFERENCES.md}
bagian~5, termasuk dua koreksi judul/tahun terhadap teks kode.
